\documentclass[10pt,a4paper]{amsart}
\usepackage[margin=19mm]{geometry}
\usepackage{amsmath,amssymb,mathtools}
\usepackage[scr=rsfs,cal=euler]{mathalfa}
\usepackage{xfrac}
\usepackage[unicode,hidelinks,bookmarks=false]{hyperref}
\newcommand{\ie}{i.e.\ }
\newcommand{\cf}{cf.\ }
\newcommand{\resp}{respectively\ }
\newcommand{\Subsection}{\subsection}
\providecommand{\nobreakdash}{\nobreak\hbox{-}}
\setlength{\emergencystretch}{2em}
\multlinegap0pt
\usepackage{amssymb}
\usepackage{tikz}
\usepackage{algorithm}\usepackage{algpseudocode}
\newtheorem{lemma}{Lemma}
\newcommand{\lcm}{lcm}
\newcommand{\ord}{ord}
\title{Complex Representations of Groups and Involutions of its Automorphisms}
\author{Venkata Subbaiah Yerrapati, Rahul Dixit, Ajay Kumar Shukla}
\date{Source: arXiv:2605.23195v1 (CC BY 4.0)}
\begin{document}
\maketitle

\noindent\emph{Source and licence.} Complex Representations of Groups and Involutions of its Automorphisms. Version arXiv:2605.23195v1 (CC BY 4.0), distributed by arXiv under the Creative Commons Attribution 4.0 International licence, http://creativecommons.org/licenses/by/4.0/, as recorded in the arXiv licence metadata for this exact version and shown on the abstract page https://arxiv.org/abs/2605.23195v1. Full licence notice: \texttt{licenses/arxiv-2605.23195-CC-BY-4.0.txt}.\par\smallskip

\noindent\textit{Editorial scope.} Counted unit: the article's lemma all\_involutions\_of\_greater\_order\_3 (source0.tex lines 305--307). The article's complete original proof of it is delivered with it (source0.tex lines 309--412, one \texttt{proof} environment of 104 lines). No mathematics is written, repaired or re-derived: the statement and the proof are the article's own lines, in the article's own notation, and no retained line is substituted or re-flowed.  No further source-local context is needed: every cross-reference of every retained line resolves inside the two delivered spans.  Nothing was omitted from any retained span, so the package carries no omission record.  No retained line contains a citation, so the wrapper carries no bibliography and no bibliography entry.  No retained line contains a figure environment, an image-inclusion command or drawing code, so no image is delivered and the package declares no asset dependency.  Editorial formatting only, in addition: an AMS \texttt{amsart} preamble that restates the article's own declarations and macros used by the retained lines, namely the environments \texttt{lemma} on separate counters, together with the standard packages those lines need.  The retained environments print in this document as Lemma 1 in order of appearance, whereas the article prints the same items with the numbering of its own full document.  The complete native source of the article as distributed in the arXiv e-print for this exact version is bundled unchanged as \texttt{sources/201193/source0.tex}.
    \begin{lemma}\label{all_involutions_of_greater_order_3}
            Let $n$ be a natural number with $n \neq 6$, $x\in\mathfrak{S}_n$ such that $\ord (x) \geq 3$ and $\ord(x)$ is odd. Suppose that the automorphism corresponding to $x$ is $\sigma : \mathfrak{S}_n \rightarrow \mathfrak{S}_n$, defined by $\sigma(y) = xyx^{-1}$ for all $y \in \mathfrak{S}_n$. If $\sigma(y) = y^{-1}$, then $\ord (y) = 1$ or $\ord (y) = 2$.
        \end{lemma}

        \begin{proof}
            Suppose, on the contrary, there exist $y \in \mathfrak{S}_n$ with $\ord (y) = k \geq 3$ such that $\sigma(y) = y^{-1}$, then we explore the following possibilities.
            \begin{itemize}
                \item[Case A.] If $y$ is a product of m-disjoint cycles of same order $k$.
                \begin{align*}
                    y = (y_1, y_2, y_3, \ldots , y_k)(y_{k + 1}, y_{k + 2}, \ldots, y_{2k})\cdots(y_{(m - 1)k + 1}, y_{(m - 1)k + 2}, \ldots ,y_{mk})
                \end{align*}
                The following are the possible cases of $y$ such that $y$ maps to $y^{-1}$ under $\sigma$. 
                \begin{itemize}
                    \item[Case (1)] Each cycle of $y$ is mapping to its own inverse.
                    \begin{align*}
                        y = (y_1, y_2, \ldots, y_k)\gamma
                    \end{align*}
                    where $\gamma$ is the product of remaining cycles that are appearing in $y$. Since $(y_1, y_2, y_3, \ldots , y_k)$ maps to its inverse $(y_k, \ldots y_2, y_1) = y^{-1}$ under $\sigma$. Assuming that cycle is already rearranged in a way that implies
                    % $x$ is of the form $(y_1, y_k)(y_2, y_{k - 1})\ldots(y_{\frac{k}{2}}, y_{\frac{k}{2} + 1})\delta_1$ if $k$ is even and $(y_1, y_k)(y_2, y_{k - 1})\ldots(y_{\frac{k - 1}{2}}, y_{\frac{k - 1}{2} + 2})\delta_2$ when $k$ is odd, 
                    \begin{align*}
                        x = \begin{cases}
                            (y_1, y_k)(y_2, y_{k - 1})\cdots(y_{\frac{k}{2}}, y_{\frac{k}{2} + 1})\delta_1 & \text{if } k \text{ is even}\\
                            (y_1, y_k)(y_2, y_{k - 1})\cdots(y_{\frac{k - 1}{2}}, y_{\frac{k - 1}{2} + 2})\delta_2 & \text{ if } k \text{ is odd}
                        \end{cases}
                    \end{align*}
                    where $\delta_1, \delta_2$ are the product of disjoint transpositions which sends the cycles of $\gamma$ to their own inverses in the respective cases. So
                    \begin{align*}
                        \ord(x) &= \lcm\{2, \ord(\delta_1)\} \text{ or } lcm\{2, \ord(\delta_2)\}\\
                        &= 2j_1 \text{ for some $j_1 \in \mathbb{N}$}
                    \end{align*}
                    which contradicts the fact that $\ord(x)$ is odd.
                    \item[Case (2).] Each cycle of $y$ is mapping to other cycle's inverse.
                    \begin{align*}
                        y = (y_1, y_2, y_3, \ldots, y_k)(y_{k + 1}, y_{k + 2}, \ldots, y_{2k}) \ldots (y_{(m - 1)k + 1}, y_{(m - 1)k + 2}, \ldots, y_{mk})
                    \end{align*}
                    We prove it by induction on number of disjoint cycles. Suppose $(y_1, y_2, y_3, \ldots , y_k) \mapsto (y_{(s- 1)k + 1}, y_{(s- 1)k + 2}, \ldots , y_{sk})$ and $(y_{(s- 1)k + 1}, y_{(s- 1)k + 2}, \ldots , y_{sk}) \mapsto (y_1, y_2, y_3, \ldots , y_k)$. Assuming that cycle is already rearranged in a way, which implies that $x$ is of the form
                    \begin{align*}
                        x = (y_1, y_{sk})(y_2, y_{sk - 1}) \cdots (y_k, y_{(s - 1)k + 1})\ \delta_3
                    \end{align*}
                    Where $\delta_3$ is product of other disjoint cycles and the order of $x$ is \textit{lcm}$\{2, \ord(\delta_1)\}$  which is equal to $2j_2$ for some $j_2 \in \mathbb{N}$, which contradicts the fact that $\ord(x)$ is odd.
    
                    When $3$-cycles are mapping to their inverses cyclically.
                    \begin{align*}
                        (y_1, y_2, y_3, \ldots , y_k) &\mapsto (y_{(s- 1)k + 1}, y_{(s- 1)k + 2}, \ldots , y_{sk})^{-1}\\
                        (y_{(s- 1)k + 1}, y_{(s- 1)k + 2}, \ldots , y_{sk})&\mapsto (y_{(r- 1)k + 1}, y_{(r- 1)k + 2}, \ldots , y_{rk})^{-1}\\
                        (y_{(r- 1)k + 1}, y_{(r- 1)k + 2}, \ldots , y_{rk})&\mapsto (y_1, y_2, y_3, \ldots , y_k)^{-1}
                    \end{align*} Assuming that cycle is already rearranged in a way, which implies that $x$ is of the form
                    \begin{align*}
                        x = (y_1, y_{sk}, y_{(r - 1)k + 1}, y_k, y_{(s - 1)k + 1}y_{rk})\delta_4
                    \end{align*}
                    where $\delta_4$ is product of other disjoint cycles and the order of $x$ \textit{lcm} $\{6, \ord(\delta_4)\}$ which is equal to $2j_3$ for some $j_3 \in \mathbb{N}$ which contradicts the fact that $\ord(x)$ is odd. Now we assume it is true for $m_1$-disjoint cycles where each cycle is mapping to other cyclically.
                    \begin{align*}
                        (y_1, y_2, y_3, \ldots , y_k) &\mapsto (y_{(s- 1)k + 1}, y_{(s- 1)k + 2}, \ldots , y_{sk})^{-1}\\
                        (y_{(s- 1)k + 1}, y_{(s- 1)k + 2}, \ldots , y_{sk})&\mapsto (y_{(r- 1)k + 1}, y_{(r- 1)k + 2}, \ldots , y_{rk})^{-1}\\
                        (y_{(r- 1)k + 1}, y_{(r- 1)k + 2}, \ldots , y_{rk})&\mapsto \cdots\\
                        (y_{(m_1- 1)k + 1}, y_{(m_1- 1)k + 2} \ldots y_{m_1k})&\mapsto (y_1, y_2, y_3, \ldots , y_k)^{-1}
                    \end{align*}
                    \begin{itemize}
                        \item[Case (2.a)] If $m_1$ is odd, then on assuming that the cycle is already rearranged in a way which implies that $x$ is of the form
                        \begin{align*}
                            x = (y_1, y_{2k}, y_{2k + 1}, \ldots , y_{(m_1 - 1)k + 1}, y_k, y_{k + 1}, \ldots, y_{(m_1 - 2)k + 1, y_{m_1k}})\delta_5
                        \end{align*} where the cycle is of order $2m_1$ and $\delta_5$ is disjoint product of other cycles
                        \item[Case (2.b)] If $m_1$ is even, then on assuming that the cycle is already rearranged in a way which implies that $x$ is of the form
                        \begin{align*}
                            x = (y_1, y_{2k}, y_{2k + 1}, \ldots , y_{m_1k})\delta_6
                        \end{align*} where the cycle is of order $m_1$ and $\delta_6$ is disjoint product of other cycles
                    \end{itemize}
                    Now we prove for $m_1 + 1$ disjoint cycles where each cycle is mapping to other cyclically.
                    \begin{align*}
                        &(y_1, y_2, y_3, \ldots , y_k) \mapsto (y_{(s- 1)k + 1}, y_{(s- 1)k + 2}, \ldots , y_{sk})\mapsto \\
                        &(y_{(r- 1)k + 1}, y_{(r- 1)k + 2}, \ldots , y_{rk})\mapsto \ldots
                        (y_{(m_1- 1)k + 1}, y_{(m_1- 1)k + 2} \ldots y_{m_1k})\mapsto\\
                        & (y_{(m_1)k + 1}, y_{m_1k + 2} \ldots y_{(m_1 + 1)k})\mapsto (y_1, y_2, y_3, \ldots , y_k)
                    \end{align*}
                    \begin{itemize}
                        \item[Case (2.b.i)] If $m_1$ is odd i.e., $m_1 + 1$ is even, then we have $x$ of the form
                        \begin{align*}
                            (y_1, y_{2k}, \ldots y_{(m_1 - 1)k + 1}, y_{(m_1 + 1)k})\delta_7
                        \end{align*}where the cycle is of order $m_1 + 1$ and $\delta_6$ is disjoint product of other cycles ($m_1 + 1 = 2j_4$ for some $j_4 \in \mathbb{N}$.) and the order of $x$ is \textit{lcm}$\{2j_4, \ord(\delta_7)\}$ which is equal to $2j_5$ for some $j_5 \in \mathbb{N}$ which contradicts the fact that $\ord(x)$ is odd.
                        
                        \item[Case (2.b.ii)] If $m_1$ is even i.e., $m_1 + 1$ is odd, then we have $x$ of the form
                        \begin{align*}
                            (y_1, y_{2k}, \ldots y_{m_1k + 1}, y_{(m_1 + 1)k}, y_k, y_{k + 1}, \ldots y_{(m_1 - 1)k + 1}, y_{(m_1 + 1)k})\delta_8
                        \end{align*} where the cycle is of order $2(m_1 + 1)$ and $\delta_8$ is disjoint product of other cycles and the order of $x$ is \textit{lcm}$\{2(m_1 + 1), \ord(\delta_8)\}$ which is equal to $2j_6$ for some $j_6 \in \mathbb{N}$ which contradicts the fact that $\ord(x)$ is odd.
                    \end{itemize}
                    Hence by Mathematical Induction this is true for all natural numbers greater than or equal to 3.
                    \item Some of them mapped to their inverse and others are mapped to other's inverse. This case follows either from first part or second part of this case.
                \end{itemize}
                \item[Case B.] If $y$ is product of different order disjoint cycles. 
                Since $y$ contains $(y_1, y_2, \ldots y_k)$ and no other cycle of same order. So, $(y_1, y_2, \ldots y_k)$ should map to $(y_k, y_{k - 1}, \ldots y_1)$ under $\sigma$
                    \begin{enumerate}
                        \item[Case (1).] If $k$ is even. \begin{align*}
                            y = (y_1, y_2, \ldots y_k) \rightarrow (y_k, \ldots y_2, y_1) = y^{-1}
                        \end{align*}
                        Which means $y_1$ goes to $y_k$ and $y_k$ goes to $y_1$. A similar arrangement is assigned to other elements. This shows that $x$ is the product of $k/2$ disjoint transpositions, which implies $\ord (x)$ is even, which is a contradiction to the fact that $\ord(x)$ is odd.
                        \item[Case (2).] If $k$ is odd.
                        \begin{align*}
                            y = (y_1, y_2, \ldots y_k) \rightarrow (y_k, \ldots y_2, y_1) = y^{-1}
                        \end{align*}
                        Which means $y_1$ goes to $y_k$ and $y_k$ goes to $y_1$. A similar arrangement is assigned to other elements. This shows that $x$ is the product of $(k-1)/2$ disjoint transpositions, which implies $\ord(x)$ is even order, which is again a contradiction to the fact that $\ord(x)$ is odd.
                    \end{enumerate}
                \item[Case C.] If $y$ is a product of disjoint cycles such that some of them have the same order and some have different orders.

                This is followed by similar arguments of Case (B).
            \end{itemize}
            
            This completes the proof.
        \end{proof}

\end{document}
